%% ma: L12-B05-0018
%% lop: 12
%% chuong: 1
%% bai: 5
%% dang: 12.5.1
%% loai: TLN
%% muc: VD
%% dap_an: "75"
%% nguon: "Ảnh giáo viên gửi 11/10/2026 (Đề kiểm tra giữa HK1 Toán 12, THPT Việt Đức 2026-2027, mã đề 0103), Phần III Câu 4"
%% kien_thuc: "Bài 5 (bài toán tối ưu lợi nhuận bằng đạo hàm)"
%% trang_thai: chinh-thuc
%% ghi_chu: "Cùng ý tưởng với các câu tối ưu lợi nhuận của dạng 12.5.1 đã duyệt."
\begin{ex}
Một cửa hàng dự kiến bán $x$ sản phẩm mỗi tháng. Giá bán mỗi sản phẩm là $p(x)=600-0{,}5x$ (nghìn đồng) và số sản phẩm mỗi tháng không vượt quá $800$ sản phẩm. Tổng chi phí sản xuất là $C(x)=100x+50000$ (nghìn đồng). Lợi nhuận lớn nhất mà cửa hàng có thể đạt được là bao nhiêu triệu đồng?
\shortans{75}
\loigiai{Lợi nhuận là $L(x)=x\,p(x)-C(x)=600x-0{,}5x^2-100x-50000=-0{,}5x^2+500x-50000$ (nghìn đồng), với $0\le x\le800$.

$L'(x)=-x+500=0\Leftrightarrow x=500\in[0;800]$. $L'(x)>0$ trên $(0;500)$ và $L'(x)<0$ trên $(500;800)$ nên $L$ lớn nhất tại $x=500$: $L(500)=75000$ nghìn đồng $=75$ triệu đồng.}
\end{ex}
